\chapter*{General Introduction}
\addcontentsline{toc}{chapter}{General Introduction}
\markboth{General Introduction}{General Introduction}

Sooner or later, every large institution has to write down who is allowed to decide what. In a development bank, that question has teeth. An approval signed by the wrong person is a governance failure, and auditors find those. \\[0.3cm]
The African Development Bank Group keeps its answer in a document called the Delegation of Authority Matrix, the DAM, which goes through the Bank's activities one by one and records who starts each task, who has to check it, who reviews it, who approves it, and who simply has to be told it happened. \\[0.3cm]
The document is accurate. Getting at what it says is the problem. The full DAM runs past two hundred and fifty pages. Nearly all of it sits inside wide tables, and because those tables carry more columns than a page can hold, the role names at the top are printed sideways. Authority is written in a shorthand of single letters, and a small digit beside a letter can change what it means. So a staff member who wants to know whether they can sign off on a mission program has to find the right row, in the right table, in the right chapter, follow a rotated header down to work out which column belongs to which role, and then decode the character where the two meet. People do manage this. They manage it slowly and not always correctly. \\[0.3cm]
None of the bank's priorities get delivered unless decisions are made correctly by people entitled to make them. Take RDGN, the regional office covering Tunisia, Algeria, Morocco, Mauritania, and Libya. Staff there spend their days on exactly the activities the Matrix governs: Organizing missions, dealing with co-financiers, and committing funds against country programs. Each needs an answer to the same four questions. who starts it, who checks it, who approves it, who needs to know. Delegation of authority is not paperwork sitting alongside the real work. It is how the real work gets done, and anything that makes it harder to consult slows the Bank down. \\[0.3cm]
The obvious thing is to try is a chatbot. It does not work here, and the reason is worth being precise aboput. Language models write fluently whether or not they are right, and fluent text about who controls financial approval is dangerous when it is wrong. A model that invents a plausible-sounding role, or quietly leaves out one of the three people who should have been informed, hands you something that reads exactly like correct answer. For a compliance document, that is worse than no answer. \\[0.3cm]
So the system described here is built around that problem rather than hoping to avoid it. The DAM is parsed into a structured knowledge base. That knowledge base becomes a graph of activities and the roles responsible for them, and questions are answered by pulling real records out of the graph. There is a language model involved, but it never decides what the facts are, It is handed facts that have already been retrieved and asked to put them into a sentence, and before anyone sees that sentence it is checked, word by word, against the facts it came from. If a role name. If a role name is missing or changed, the sentence is thrown away and a plain one is shown instead. \\[0.5cm]
The report runs to five chapters:\\ [0.3cm]
\textbf{Chapter 1} sets the scene: the host instituion, the client, how the Matrix is consulted today, where that breaks down, what existing tools offer, and the requirements and architecture that follow from all of it\\[0.3cm]
\textbf{Chapter 2} looks hard at the source document. there is no tidy dataset here, only a PDF whose structure lives in its layout, so this chapter works through how that layout carries meaning and where automated extraction goes wrong on it. \\[0.3cm]
\textbf{Chapter 3} is the pipeline that turns the document into data, stage by stage, along with the checks applied at each one and the knowledge graph built on the output\\[0.3cm]
\textbf{Chapter 4} covers the agent: how questions finds the right activity, how an answer is assembled, what language model allowed to do, and the check that keeps it honest \\[0.3cm]
\textbf{Chapter 5} presents what was actually delivered, including the features added in the project's second phase, the analytics dashboard, the testing behind it, and how it is deployed. The general conclusion then takes stock of what worked, what did not, and what should be tackled next
